% Source: https://artofproblemsolving.com/community/c3h2960617_the_millionth_user
% Stable identity: aops-topic-2960617-q23
Oppositely marked winning lines cannot intersect. The only disjoint winning lines are two distinct rows or two distinct columns. Count row configurations first. Choose a triangle row and a different circle row in $3\cdot2=6$ ways, then fill the remaining row in $8$ ways, initially giving $48$. Configurations with three monochromatic rows of both types are counted twice: there are $2^3-2=6$. Hence there are $42$ row configurations. Similarly there are $42$ column configurations. A row configuration cannot also be a column configuration, since a monochromatic row forbids an oppositely marked column. Thus the answer is $\boxed{84}$.
\begin{asy}
usepackage("amsmath");
size(7cm);
for(int i=0;i<=3;++i){draw((i,0)--(i,3));draw((0,i)--(3,i));}for(int i=0;i<3;++i){label("$\triangle$",(i+.5,2.5));label("$\bigcirc$",(i+.5,.5));}label("opposite winning rows",(1.5,-.5));
\end{asy}
